% 用户上传的 Irodov Arihant 底本，PDF 页88--89，原书印刷页79--80。
% 原题 146--159；原书图2.4两种T-S循环已经对照图像检查。
\thermoentropyproblem{146}{79}
一摩尔定容摩尔热容为 $C_V$ 的理想气体经历一个过程，
在过程中其熵与温度满足 $S=\alpha/T$，
其中 $\alpha$ 为常数。气体温度由 $T_1$ 变化至 $T_2$。
求：(a) 过程摩尔热容 $C(T)$；
(b) 传入气体的热量 $Q$；
(c) 气体做功 $A$。
\exsolution
过程摩尔热容的定义为
$C(T)=T\,\mathrm dS/\mathrm dT$。
从 $S=\alpha/T$ 求导得
\[
\boxed{C(T)=-\frac{\alpha}{T}}.
\]
沿准静态过程
$\delta Q=T\,\mathrm dS=-\alpha\,\mathrm dT/T$，
故
\[
\boxed{Q=-\alpha\int_{T_1}^{T_2}\frac{\mathrm dT}{T}
=\alpha\ln\frac{T_1}{T_2}}.
\]
一摩尔理想气体的内能变化
$\Delta U=C_V(T_2-T_1)$，
热力学第一定律 $Q=\Delta U+A$ 给出
\[
\boxed{A=\alpha\ln\frac{T_1}{T_2}
+C_V(T_1-T_2)}.
\]
\excheck $C(T)$ 是\emph{过程}热容而非 $C_V$；
只要 $T>0$，其符号取决于 $\alpha$，
热量和外功都必须保留代数符号。

\thermoentropyproblem{147}{79}
某工质在图2.4(a)、(b) 所示的两种
$T$--$S$ 三角形循环中工作，
循环最高温与最低温之比为
$n=T_h/T_c>1$。求两种循环的热效率。
\exsolution
在温度—熵图上，沿可逆支路的热量
$\delta Q=T\,\mathrm dS$；
竖直支路熵恒定，不交换热。
设两端水平熵间隔为 $\Delta S>0$。
(a) 按原图箭头，顶部恒温支路向右，
吸热
$Q_{\rm in}=T_h\Delta S$；
倾斜支路返回时向外放热，
由于直线在 $T$--$S$ 图上的面积为梯形，
放热模为
$Q_{\rm out}=(T_h+T_c)\Delta S/2$。
所以
\[
\boxed{\eta_a=1-\frac{Q_{\rm out}}{Q_{\rm in}}
=1-\frac{T_h+T_c}{2T_h}
=\frac{n-1}{2n}}.
\]
(b) 原图沿上升斜线向右吸热，
$Q_{\rm in}=(T_h+T_c)\Delta S/2$；
底部等温线向左排热
$Q_{\rm out}=T_c\Delta S$，
因此
\[
\boxed{\eta_b=1-\frac{2T_c}{T_h+T_c}
=\frac{n-1}{n+1}}.
\]
\excheck 两循环都是顺时针正功方向；
必须依据原图箭头判别吸热段，
而不能只看两个三角形的外形。
两个效率均小于卡诺上界 $1-1/n$。

\thermoentropyproblem{148}{79}
两个彼此绝热的容器用带阀门的短管相连，
其中一个装有 $\nu=2.2$ mol 理想气体，
另一个抽成真空。阀门打开后，
气体的可及体积增大为原来的 $n=3.0$ 倍。
求气体的熵增。
\exsolution
将两容器整体视为隔热、刚性系统。
气体向真空自由膨胀，没有外功也没有热交换，
因此 $\Delta U=0$。
理想气体内能只依赖温度，
故末态温度等于初温。
熵为状态函数，可以用假想可逆等温膨胀
连接相同初末态：
\[
\Delta S=\int_{V_1}^{V_2}\frac{\nu R}{V}\,\mathrm dV
=\nu R\ln(V_2/V_1).
\]
于是
\[
\boxed{\Delta S=\nu R\ln n
=(2.2)(8.314)\ln3
\simeq20.1\,\mathrm{J/K}}.
\]
\excheck 实际自由膨胀不可逆，不能沿实际路径写
$\delta Q_{\rm actual}=T\,\mathrm dS$；
熵积分使用的是假想可逆路径。

\thermoentropyproblem{149}{79}
一个绝热圆筒被无质量活塞分成体积相等的两半。
左半部装有一摩尔绝热指数为 $\gamma$、
初温 $T_0$ 的理想气体，右半部是真空。
释放活塞，使气体充满整个圆筒。
随后缓慢将活塞移回原位置。
求这两个过程完成之后气体的内能增量和熵增。
\exsolution
释放时气体向真空自由膨胀到原体积的两倍；
由于整体绝热、无外功，
$\Delta U_1=0$ 且 $T_1=T_0$。
自由膨胀产生熵
$\Delta S_1=R\ln2$。
第二步缓慢移动活塞，
把气体由 $2V_0$ 压回 $V_0$。
按题设圆筒隔热与缓慢准静态运动，
这是可逆绝热压缩；
由 $TV^{\gamma-1}=\text{常数}$，
\[
T_f=T_0\left(\frac{2V_0}{V_0}\right)^{\gamma-1}
=2^{\gamma-1}T_0.
\]
一摩尔总内能增加
\[
\boxed{\Delta U=C_V(T_f-T_0)
=\frac{RT_0}{\gamma-1}(2^{\gamma-1}-1)}.
\]
第二步熵不变，故总熵增为
\[
\boxed{\Delta S=R\ln2}.
\]
\excheck 原书答案文本抽取可能把
$2^{\gamma-1}$ 显示成“$2\gamma-1$”；
由绝热温度体积关系确认指数结构。

\thermoentropyproblem{150}{79}
同一初始态的理想气体在不与外界交换热量的条件下
膨胀到相同的终体积 $V$。
比较：(a) 快速膨胀；(b) 极其缓慢膨胀。
末态压强是否相同？
\exsolution
对同一初态、相同终体积，
如果慢膨胀始终可逆绝热，
末态满足 $pV^\gamma=\text{常数}$。
而实际快速膨胀通常产生不可逆耗散，
同一终体积下其熵高于可逆绝热终态。
理想气体的熵差可写为
\[
\Delta S=C_V\nu\ln(T_f/T_i)
+\nu R\ln(V_f/V_i).
\]
两过程的初态与终体积相同，
较高的终态熵对应较高的终态温度。
由 $p_f=\nu RT_f/V_f$，
得出
\[
\boxed{p_{f,\rm fast}>p_{f,\rm slow}}
\]
（快速过程出现熵产生、而慢过程可逆时）。
\excheck 这是一项\emph{关于实际快过程与可逆慢过程}
的比较：如果两个过程具有相同的实际功和能量耗散，
就不能仅凭“快慢”二字确定严格的压强差。

\thermoentropyproblem{151}{80}
一个绝热容器被隔板分成两部分，
大的一侧体积是小的一侧的 $n=2.0$ 倍。
小室内有 $\nu_1=0.30$ mol 氮气，
大室内有 $\nu_2=0.70$ mol 氧气。
两气体初温相同。
在隔板上穿孔让两气体混合，
视为理想气体，求系统熵增。
\exsolution
设小室容积为 $V$、大室为 $nV$。
无热交换、无体积功且两气体初温相同，
理想气体内能守恒表明末温不变。
每种气体可等温地从各自初始容积
扩散到共同总容积 $(n+1)V$。
因此
\[
\Delta S_1=\nu_1R\ln\frac{(n+1)V}{V},\qquad
\Delta S_2=\nu_2R\ln\frac{(n+1)V}{nV}.
\]
相加：
\[
\boxed{\Delta S=
\nu_1R\ln(1+n)
+\nu_2R\ln(1+1/n)}.
\]
取给定数值，
\[
\boxed{\Delta S=R(0.30\ln3+0.70\ln1.5)
\simeq5.10\,\mathrm{J/K}}.
\]
\excheck 氮、氧是不同分子，
均可占据最终整个容积；
不能将混合熵误写成同一种气体移除隔板的零熵变。

\thermoentropyproblem{152}{80}
把质量 $m_1=300\,\mathrm g$、
温度 $t_1=97^\circ\mathrm C$ 的铜块，
投入质量 $m_2=100\,\mathrm g$、
温度 $t_2=7^\circ\mathrm C$ 的水中。
假定热量计热容可忽略，整体隔热。
求达到温度平衡时系统的总熵增。
\exsolution
取铜和水的质量比热分别为
$c_1,c_2$，在本温区视为常数。
能量守恒：
\[
m_1c_1(T_f-T_1)+m_2c_2(T_f-T_2)=0,
\quad
T_f=\frac{m_1c_1T_1+m_2c_2T_2}
{m_1c_1+m_2c_2}.
\]
各自熵变由可逆积分求得：
\[
\boxed{\Delta S=
m_1c_1\ln\frac{T_f}{T_1}
+m_2c_2\ln\frac{T_f}{T_2}}.
\]
用 $c_1\simeq0.39\,\mathrm{J/(g\,K)}$、
$c_2\simeq4.18\,\mathrm{J/(g\,K)}$，
$T_1=370.15\,\mathrm K$、
$T_2=280.15\,\mathrm K$，
先算得 $T_f\simeq300\,\mathrm K$，
再代熵式：
\[
\boxed{\Delta S\simeq3.7\,\mathrm{J/K}}
\quad(c_{\rm Cu}=0.39,\ c_w=4.18\ \mathrm{J/(g\,K)}).
\]
\excheck 熵变应使用\emph{开尔文温度}取对数，
不能将摄氏温度比值直接代入。
本版书末给约 $4.4\,\mathrm{J/K}$；
该数值依赖所用铜、水比热表和温度积分近似，
本稿明确列出所用常数与公式，不直接照搬近似数值。

\thermoentropyproblem{153}{80}
两只相同体积、彼此绝热的刚性容器，
各装一摩尔\emph{同一种}理想气体，
初温分别为 $T_1,T_2$。
其定容摩尔热容为 $C_V$。
打开阀门后两容器气体达到共同平衡，
求气体总熵增并证明非负。
\exsolution
整体绝热且体积不变，
相同气体每容器各一摩尔，
故能量守恒给
$C_V(T_1+T_2)=2C_VT_f$，
\[
T_f=\frac{T_1+T_2}{2}.
\]
按各自一摩尔的熵变加和；
两容器总体积相同、最终混合的气体为同一种，
若用分室可逆路径计算，体积项总体抵消，
因而
\[
\Delta S=C_V\ln(T_f/T_1)
+C_V\ln(T_f/T_2)
=\boxed{C_V\ln\frac{(T_1+T_2)^2}{4T_1T_2}}.
\]
由算术—几何平均不等式
$(T_1+T_2)^2\ge4T_1T_2$，
所以
\[
\boxed{\Delta S\ge0},
\]
且仅在 $T_1=T_2$ 时等号成立。
\excheck 两种气体若是不同种类，还需包括混合熵；
本题明确为\emph{同一种气体}，不可加虚假混合熵。

\thermoentropyproblem{154}{80}
有 $N$ 个氦原子处于室温、体积为
$1.0\,\mathrm{cm^3}$ 的立方容器中。
求：(a) 所有原子在同一瞬间全部聚集于一半容积的概率；
(b) 若要求这一事件在约
$t=10^{10}$ 年内发生一次，需要大致多少个原子？
\exsolution
若原子位置独立且在两半室等概率分布，
指定某一半的每个原子都位于其中的概率为 $1/2$，
故
\[
\boxed{P=2^{-N}}.
\]
(b) 设重新抽样出近似独立位置分布的时间尺度为
$\tau$，若观察时长为 $t$，
估计独立尝试次数为 $K\simeq t/\tau$。
每次概率为 $2^{-N}$，
平均发生次数 $K2^{-N}\sim1$，故
\[
\boxed{N\simeq\frac{\ln(t/\tau)}{\ln2}}.
\]
取室温原子通过厘米量级容器的时间
$\tau\simeq10^{-5}\,\mathrm s$、
$t\simeq3.16\times10^{17}\,\mathrm s$，
得到
$\boxed{N\simeq75\ \text{个，约80个}}$。
\excheck (b) 只是基于近似独立尝试的数量级估计；
真实微观位置关联时间不严格等于一次过箱时间。
(a) 按原书是\emph{指定}的一半区域，
若接受左右任一半，概率须改为 $2^{1-N}$。

\thermoentropyproblem{155}{80}
把 $N=10$ 个相同分子分配到圆筒等体积的两半，
求最概然分配的统计权重及其发生概率。
\exsolution
各分子可独立位于左半或右半，
宏观态“左半有 $k$ 个”对应的微观排列数为
\[
\Omega_k=\binom Nk=\frac{N!}{k!(N-k)!}.
\]
对于偶数 $N=10$，最概然数为 $k=N/2=5$，
所以
\[
\boxed{\Omega_{\rm max}=\binom{10}{5}
=\frac{10!}{(5!)^2}=252}.
\]
共有 $2^{10}$ 个等概率左右分配，
概率为
\[
\boxed{P_{\rm max}=\frac{252}{2^{10}}
=\frac{252}{1024}\simeq24.6\%}.
\]
\excheck “分子相同”指它们物理种类相同；
排列计数仍可以对可追踪的微观粒子位置计数，
不应把所有分配只算成一个状态。

\thermoentropyproblem{156}{80}
容器装有 $N$ 个理想气体分子，
把容积设想划为相等的 $A,B$ 两半。
求 $A$ 半中有恰好 $n$ 个分子的概率。
特别列出 $N=5$ 时
$n=0,1,2,3,4,5$ 的概率。
\exsolution
每个分子进入 $A$ 的概率为 $1/2$，
所有分子位置独立时，
恰好 $n$ 个在 $A$ 的概率服从二项分布：
\[
\boxed{P_n=\binom Nn
\left(\frac12\right)^n
\left(\frac12\right)^{N-n}
=\frac{N!}{n!(N-n)!\,2^N}}.
\]
$N=5$ 的六个概率依次为
\[
\boxed{
P_0=\frac1{32},\
P_1=\frac5{32},\
P_2=\frac{10}{32},\
P_3=\frac{10}{32},\
P_4=\frac5{32},\
P_5=\frac1{32}}.
\]
\excheck 六项相加
$(1+5+10+10+5+1)/32=1$，
满足概率归一化。

\thermoentropyproblem{157}{80}
体积为 $V_0$ 的容器内有 $N$ 个理想气体分子。
任选其中体积为 $V$ 的固定子区域，
求其中恰有 $n$ 个分子的概率，
并讨论 $V=V_0/2$。
\exsolution
位置均匀分布时，单个分子落入子区域的概率
$q=V/V_0$。
恰有 $n$ 个分子落入的组合数为 $\binom Nn$，
因此
\[
\boxed{P_n=\binom Nn
\left(\frac V{V_0}\right)^n
\left(1-\frac V{V_0}\right)^{N-n}}.
\]
若 $V=V_0/2$，
则 $q=1/2$，
\[
\boxed{P_n=\binom Nn2^{-N}}.
\]
\excheck 需要 $0\le V\le V_0$ 与 $0\le n\le N$；
所有 $n$ 的概率之和由二项式定理等于1。

\thermoentropyproblem{158}{80}
标准状况下的理想气体，
求一个球形取样区域的直径 $d$，
使区域内分子数的相对涨落为
$\eta=1.0\times10^{-3}$。
同时求该区域内平均分子数。
\exsolution
对均匀稀薄气体中的固定小取样区域，
当整体系统远大于该区域，
分子数涨落近似泊松分布：
$\operatorname{Var}N=\bar N$。
相对标准差为
\[
\eta=\frac{\sqrt{\operatorname{Var}N}}{\bar N}
=\frac1{\sqrt{\bar N}},
\]
故
\[
\boxed{\bar N=\eta^{-2}=10^6}.
\]
标准状况数密度
$n_0=p_0/(k_BT_0)
\simeq2.69\times10^{25}\,\mathrm{m^{-3}}$。
球体积
$\frac{\pi d^3}{6}=\bar N/n_0$，
所以
\[
\boxed{d=\left(\frac6{\pi n_0\eta^2}\right)^{1/3}
\simeq4.1\times10^{-7}\,\mathrm m
=0.41\,\mathrm{\mu m}}.
\]
\excheck 由体积求直径必须取\emph{三次方根}；
PDF 提取文本有时丢掉此根号结构。

\thermoentropyproblem{159}{80}
一摩尔单原子理想气体装在等容容器中，
初温 $T_0=300\,\mathrm K$，
加热 $\Delta T=1.0\,\mathrm K$。
求气体统计权重 $\Omega$ 变化的倍率。
\exsolution
玻尔兹曼熵公式
$S=k_B\ln\Omega$，
故统计权重比
$\Omega_f/\Omega_i=e^{\Delta S/k_B}$。
等容加热一摩尔单原子气体的熵增：
\[
\Delta S=C_V\ln\frac{T_0+\Delta T}{T_0},
\quad C_V=\frac32R=\frac32N_Ak_B.
\]
因此
\[
\boxed{\frac{\Omega_f}{\Omega_i}
=\left(1+\frac{\Delta T}{T_0}\right)^{3N_A/2}}.
\]
用十进制数量级表示：
\[
\log_{10}\frac{\Omega_f}{\Omega_i}
=\frac{3N_A}{2}\log_{10}\left(\frac{301}{300}\right)
\simeq1.31\times10^{21},
\]
即
\[
\boxed{\Omega_f/\Omega_i
\simeq10^{\,1.31\times10^{21}}}.
\]
\excheck 这是极巨大的\emph{无量纲倍率}，
不能误写为简单的 $10^{21}$。
